\Needspace{12\baselineskip}
\originalchapter{分子动力学与输运}
\label{chap:md}
\originalsection{运动方程和观测量}
在势能面 $U(\mathbf R)$ 上，经典原子满足 $m_i\ddot{\mathbf R}_i=-\nabla_iU$。轨迹比静态最小化提供更多信息：振动、扩散、时间相关和非平衡响应都依赖它。但电子激发、核量子效应或反应势的适用范围仍需单独检查。

速度Verlet把一步拆为半步速度、整步位置和再半步速度：$\mathbf v_{n+1/2}=\mathbf v_n+\Delta t\mathbf F_n/(2m)$，$\mathbf R_{n+1}=\mathbf R_n+\Delta t\mathbf v_{n+1/2}$，$\mathbf v_{n+1}=\mathbf v_{n+1/2}+\Delta t\mathbf F_{n+1}/(2m)$。对时间无显式依赖的光滑保守力，它是时间可逆的二阶算法。能量的有界小振荡与持续系统漂移应区别；增大步长可能激发最高频模而失稳。

微正则生产阶段不应持续缩放速度。恒温平衡阶段可以使用合适恒温器，但恒温器会改变时间相关，输运性质可能受其影响。固定总动量后，可运动自由度从 $3N$ 减为 $3N-3$，温度估计 $T=2K/[(3N-3)\kb]$。温度沿一条有限轨迹涨落并不表示能量不守恒。

\originalsection{截断、压力和结构}
普通势能截断 $v(r)\mathbf1_{r<r_c}$ 在边界跳变；仅做势移位仍使力在边界跳变。力移位形式
\[
v_{fs}(r)=v(r)-v(r_c)-(r-r_c)v'(r_c),\quad r<r_c
\]
使势和一阶导在边界连续，改善短轨迹能量检查，但也改变了模型本身。若用它，报告中必须写明，不能把其物性当成原无限程LJ势的结果。

各向同性体系的压力可由动量通量和virial得到 $P=N\kb T/V+\sum_{i<j}\mathbf r_{ij}\cdot\mathbf F_{ij}/(3V)$。周期体系要用一致的镜像位移；具有特殊约束或长程电荷时，还会有额外项。径向分布函数 $g(r)$ 将壳层平均计数除以 $4\pi r^2\Delta r$、体相密度和参考粒子数。对有限小胞，靠近 $L/2$ 的统计与有限尺寸效应需要注意。

\originalsection{扩散和时间相关}
扩散系数由长期展开坐标的均方位移 $D=\lim_{t\to\infty}\avg{|\mathbf R(t)-\mathbf R(0)|^2}/(6t)$。若每步把位置折回盒内后直接求差，穿过边界的原子会被误判为一次大跳或被限制在盒内。应保存未折回坐标或累计穿界次数。

位移是速度积分，平稳条件下 $\avg{\Delta R^2}=2\int_0^t(t-s)\avg{\mathbf v(s)\cdot\mathbf v(0)}\dd s$，长时间得 $D=\tfrac13\int_0^\infty\avg{\mathbf v(s)\cdot\mathbf v(0)}\dd s$，即Green--Kubo关系。短时间弹道区 $\avg{\Delta R^2}\propto t^2$，晶体受限振动可形成平台；两者都不应拟合成长期扩散直线。

附带程序运行约化单位LJ体系，保存完整参数、轨迹统计和两种步长的能量波动。它核验保守力、积分和周期边界的实现。模拟既没有采用手册的Ni-EAM参数，也没有完整运行Kr两相共存熔点任务，不把短时间体系运动误称为原作业的真实材料结果。

\originalsection{系综与恒温器的物理影响}
微正则系综固定 $N,V,E$，正则系综固定 $N,V,T$，等温等压又允许体积涨落。一条有限轨迹能否代表系综还需遍历与足够混合；即使目标分布理论正确，短运行仍可能只停在一个相态。

Langevin恒温方程 $m\dot v=F-\gamma v+\sqrt{2\gamma\kb T}\eta$ 中耗散与噪声配对确保正则稳态。改变 $\gamma$ 影响动力学相关和扩散。简单每步把动能缩到固定温度，压掉正常涨落，一般不产生正确正则能量分布。Nos\'e--Hoover类确定性恒温器有扩展变量，某些小体系仍可不遍历；知道名称不能替代对实际分布的检验。

周期胞先设置密度与构型、消去总动量，经过有明确恒温方案的热化，再切换生产采样。如果要测真实Hamilton动力学输运，通常需要尽量减少恒温器对关注相关的干扰。应分别画热化与生产段，不能把系统起始的结构重排混进稳态平均。

\originalsection{步长、频率和能量检查}
对简谐模式 $\ddot q+\omega^2q=0$，Verlet稳定要求 $\Delta t\,\omega<2$，但稳定不等于足够准确。复杂体系最高振动频率或硬碰撞决定限制。光滑、时间可逆积分常保留一个附近的修正Hamilton量，能量可有有界振荡；长期漂移则可能来自不连续截断、邻居表失效、迭代力误差或步长过大。

邻居表皮肤宽度应超过两次更新之间的最大相对位移，或根据位移触发重建。不同输入种子会改变轨迹，因此比较步长应从完全同一初态开始、固定总物理时间；否则轨迹不同不能归因于步长。附带108粒子、密度 $.8$ 的力移位LJ比较使用相同初始构型和速度，$\Delta t=.002,.001$ 的总能标准差约 $.000311,.0000777$，接近二阶缩放；这只是当前短时模型的守恒证据。

\originalsection{结构函数、输运与统计误差}
$g(r)$ 描述空间两点结构，静态结构因子 $S(\mathbf q)=N^{-1}\avg{|\sum_j\e^{i\mathbf q\cdot\mathbf R_j}|^2}$ 描述波数空间涨落，二者在均匀大体系中联系为 $S(q)=1+\rho\int[g(r)-1]\e^{i\mathbf q\cdot\mathbf r}\dd^3r$。有限胞允许波数离散，最小非零波数 $2\pi/L$ 限制长波结构。

输运相关的积分需要选择物理平台窗口：太早保留弹道或振动，太晚噪声累计，短小胞还会周期回声。多个时间原点彼此重叠，不能当独立路径；以块或独立运行估误差。不同盒长、轨迹长度和温控强度的比较各解决一个不同问题。

黏度的Green--Kubo关系用平衡应力自相关 $\eta=V/(\kb T)\int_0^\infty\avg{P_{xy}(t)P_{xy}(0)}\dd t$，需要正确应力定义与长时间统计。本卷没有凭短LJ能量轨迹虚构已经收敛的黏度；此式说明如何把动力学和宏观材料响应连接。
